% ============================================================================
% SECCION 5: CONCLUSIONES
% ============================================================================

\section{Conclusiones}
\label{sec:conclusiones}

% ═══════════════════════════════════════════════════════════════════════════
\subsection{Hallazgos principales}
\label{subsec:hallazgos_principales}
% ═══════════════════════════════════════════════════════════════════════════

Este estudio presenta una evaluaci\'on sistem\'atica y a gran escala de la precisi\'on
estructural de AlphaFold2, basada en la comparaci\'on de 10.432 prote\'inas \'unicas
con sus correspondientes estructuras experimentales del PDB, utilizando mapeo SIFTS
para la correspondencia correcta de residuos. Los hallazgos principales se resumen
a continuaci\'on.

\subsubsection{Precisi\'on global de AlphaFold2}

\begin{enumerate}
    \item \textbf{RMSD mediano de \SI{1,26}{\angstrom} para un \textit{dataset} diversificado}:
    Sobre 10.432 prote\'inas con identificadores UniProt \'unicos, el RMSD mediano se
    sit\'ua dentro del rango de variabilidad experimental t\'ipico
    (\SIrange{0,5}{2,0}{\angstrom}). La media aritm\'etica (\SI{3,56}{\angstrom}) es
    sustancialmente mayor debido a la distribuci\'on log-normal con cola derecha
    pronunciada.

    \item \textbf{65,3\% de predicciones con RMSD $< \SI{2}{\angstrom}$}: Aproximadamente
    dos tercios de las prote\'inas analizadas presentan una concordancia con la estructura
    experimental que las hace \'utiles para la mayor\'ia de aplicaciones
    bioinform\'aticas y de dise\~no molecular.

    \item \textbf{41,5\% con RMSD $< \SI{1}{\angstrom}$}: M\'as de cuatro de cada diez
    predicciones alcanzan precisi\'on sub-angstrom, un nivel donde la predicci\'on
    computacional es indistinguible de la determinaci\'on experimental dentro del
    margen de error cristalográfico.

    \item \textbf{RMSD mediano de \SI{0,64}{\angstrom} para prote\'inas bien estudiadas}:
    El \textit{dataset} validado (91 prote\'inas, 5.115 pares con mapeo SIFTS expl\'icito) demuestra
    que, para prote\'inas con abundante informaci\'on evolutiva y estructural, la
    precisi\'on de AlphaFold2 es excepcional.
\end{enumerate}

\subsubsection{Impacto del mapeo de residuos}

\begin{enumerate}
    \setcounter{enumi}{4}
    \item \textbf{El mapeo SIFTS es un requisito metodol\'ogico, no una opci\'on}:
    Sin mapeo posicional expl\'icito (\texttt{PDB\_BEG} $\rightarrow$ \texttt{SP\_BEG}),
    el RMSD mediano se infla artificialmente de \SI{0,64}{\angstrom} a
    \SI{22,54}{\angstrom} ---una amplificaci\'on de $35\times$. Este artefacto se origina
    por dos causas independientes: desfase de numeraci\'on entre PDB y UniProt
    (ej: p\'eptido se\~nal de la lisozima, desfase de 18 posiciones) y selecci\'on
    incorrecta de cadena en PDB multim\'ericos.

    \item \textbf{No existe un ``sweet spot'' de tama\~no}:
    El aparente rango \'optimo en prote\'inas de 250--300 residuos, identificado en
    versiones anteriores de este estudio, era un artefacto del desfase de numeraci\'on.
    Con mapeo correcto, la precisi\'on de AlphaFold2 es uniformemente alta en todos
    los rangos de tama\~no.
\end{enumerate}

\subsubsection{An\'alisis jer\'arquico}

\begin{enumerate}
    \setcounter{enumi}{6}
    \item \textbf{Gradiente de precisi\'on qu\'imicamente interpretable}:
    El an\'alisis en cuatro niveles (longitud $\to$ amino\'acido $\to$ grupo funcional
    $\to$ elemento at\'omico) revela que la precisi\'on de AlphaFold2 correlaciona con
    la rigidez conformacional: los residuos del n\'ucleo hidrofóbico (alif\'aticos,
    arom\'aticos) se predicen mejor que los de superficie (cargados, polares), y los
    \'atomos de carbono del esqueleto mejor que los \'atomos de ox\'igeno y nitr\'ogeno
    de cadenas laterales flexibles.

    \item \textbf{Los \textit{outliers} tienen explicaci\'on biol\'ogica o metodol\'ogica}:
    El 9,5\% de pares con RMSD $> \SI{10}{\angstrom}$ se explica por cambios
    conformacionales ligando-inducidos (apo vs.\ holo, $\sim$35\%), dominios m\'oviles
    ($\sim$25\%), regiones intr\'insecamente desordenadas ($\sim$20\%) y errores de
    mapeo SIFTS por valores nulos en \texttt{PDB\_BEG} ($\sim$20\%). Ninguna de
    estas categor\'ias representa un fallo intr\'inseco de AlphaFold2.
\end{enumerate}

% ═══════════════════════════════════════════════════════════════════════════
\subsection{Contribuciones metodol\'ogicas}
\label{subsec:contribuciones_metodologicas}
% ═══════════════════════════════════════════════════════════════════════════

M\'as all\'a de los resultados espec\'ificos sobre AlphaFold2, este estudio aporta
las siguientes contribuciones metodol\'ogicas:

\begin{enumerate}
    \item \textbf{Pipeline reproducible con mapeo SIFTS}: Un \textit{pipeline} completo
    (descarga $\to$ mapeo $\to$ superposici\'on $\to$ an\'alisis) implementado en Python
    con c\'odigo abierto, que puede reutilizarse para evaluar AlphaFold3 u otros
    m\'etodos de predicci\'on estructural.

    \item \textbf{Demostraci\'on del impacto del mapeo incorrecto}: La cuantificaci\'on
    expl\'icita del artefacto de $35\times$ amplificaci\'on de RMSD sirve como advertencia
    y referencia para la comunidad.

    \item \textbf{An\'alisis jer\'arquico de cuatro niveles}: La descomposici\'on de la
    precisi\'on en factores qu\'imicamente interpretables (longitud, amino\'acido, grupo
    funcional, elemento at\'omico) es una contribuci\'on novedosa que conecta
    rendimiento computacional con propiedades moleculares.

    \item \textbf{Marco bayesiano para evaluaci\'on de calidad}: El modelo Beta-Binomial
    con intervalos HDI al 95\% proporciona cuantificaci\'on de incertidumbre superior
    a las pruebas de hip\'otesis frecuentistas convencionales.

    \item \textbf{Estrategia de diversificaci\'on}: La selecci\'on de un PDB representativo
    por UniProt ID, priorizando cobertura de secuencia, es una estrategia generalizable
    para eliminar el sesgo de sobrerepresentaci\'on en bases de datos estructurales.
\end{enumerate}

% ═══════════════════════════════════════════════════════════════════════════
\subsection{Recomendaciones pr\'acticas}
\label{subsec:recomendaciones_practicas}
% ═══════════════════════════════════════════════════════════════════════════

Con base en los resultados de este estudio, se formulan las siguientes recomendaciones
para investigadores que utilicen modelos de AlphaFold2:

\begin{enumerate}
    \item \textbf{Utilizar modelos de la AFDB con confianza para la mayor\'ia de
    aplicaciones}: Con RMSD mediano de \SI{1,26}{\angstrom}, los modelos son
    suficientemente precisos para docking molecular, dise\~no racional de mutantes,
    an\'alisis de interfaces prote\'ina-prote\'ina, modelado por homolog\'ia y
    anotaci\'on funcional basada en estructura.

    \item \textbf{Mayor confianza en el n\'ucleo hidrofóbico}: Las coordenadas de
    residuos alif\'aticos y arom\'aticos (n\'ucleo estructural) son m\'as fiables que
    las de residuos cargados en la superficie. Para aplicaciones sensibles a la
    geometr\'ia del sitio activo (frecuentemente en el n\'ucleo), las predicciones
    son especialmente confiables.

    \item \textbf{Siempre utilizar mapeo SIFTS para comparaciones PDB--AlphaFold}:
    Cualquier estudio que compare estructuras por n\'umero de posici\'on sin mapeo
    posicional expl\'icito debe considerarse metodol\'ogicamente cuestionable.

    \item \textbf{Precauci\'on con la conformaci\'on holo}: AlphaFold2 predice la
    conformaci\'on apo. Para aplicaciones que requieran la conformaci\'on unida a
    ligando (ej: dise\~no de f\'armacos basado en estructura), complementar con
    docking flexible o buscar estructuras experimentales en el estado holo.

    \item \textbf{Verificar pLDDT pero no confiar exclusivamente en \'el}: El \textit{score} de
    confianza por residuo es \'util como filtro inicial, pero la validaci\'on contra
    estructuras experimentales proporciona informaci\'on cualitativamente superior.

    \item \textbf{Reportar estad\'isticas descriptivas completas}: Dada la distribuci\'on
    log-normal del RMSD, reportar mediana, rango intercuart\'ilico, percentiles clave
    (P5, P25, P75, P95) y fracci\'on por debajo de umbrales est\'andar
    (\SI{1}{\angstrom}, \SI{2}{\angstrom}, \SI{5}{\angstrom}), en lugar de \'unicamente
    la media aritm\'etica.

    \item \textbf{Diversificar \textit{datasets} de evaluaci\'on}: Seleccionar un PDB
    representativo por identificador UniProt para evitar conclusiones sesgadas por
    sobrerepresentaci\'on de prote\'inas muy estudiadas.
\end{enumerate}

% ═══════════════════════════════════════════════════════════════════════════
\subsection{Trabajo futuro}
\label{subsec:trabajo_futuro}
% ═══════════════════════════════════════════════════════════════════════════

Los resultados y limitaciones de este estudio sugieren las siguientes l\'ineas de
investigaci\'on futura:

\begin{enumerate}
    \item \textbf{Evaluaci\'on de AlphaFold3}: Cuando los modelos de AlphaFold3
    \citep{abramson2024accurate} est\'en disponibles para descarga masiva, repetir este
    an\'alisis con el mismo \textit{pipeline} permitir\'a una comparaci\'on directa y cuantitativa
    de la mejora entre versiones. De particular inter\'es es la capacidad de AlphaFold3
    para modelar m\'ultiples conformaciones mediante su m\'odulo de difusi\'on, lo que
    podr\'ia reducir los RMSDs asociados a diferencias apo/holo.

    \item \textbf{An\'alisis temporal}: Estratificar las 10.432 prote\'inas seg\'un la
    fecha de dep\'osito en el PDB (antes vs.\ despu\'es de abril 2018, fecha de corte
    del entrenamiento de AlphaFold2) para cuantificar la degradaci\'on de rendimiento
    en prote\'inas genuinamente novedosas.

    \item \textbf{An\'alisis por dominio}: Descomponer prote\'inas multidominio y
    evaluar cada dominio independientemente, usando clasificaciones CATH o SCOP. Esto
    proporcionar\'ia una evaluaci\'on m\'as justa de prote\'inas con bisagras flexibles
    y reducir\'ia el n\'umero de ``falsos \textit{outliers}''.

    \item \textbf{An\'alisis por resoluci\'on cristalográfica}: Estratificar por
    resoluci\'on de la estructura PDB ($< \SI{1,5}{\angstrom}$, \SIrange{1,5}{2,5}{\angstrom},
    $> \SI{2,5}{\angstrom}$) para separar el error de predicci\'on del error
    experimental.

    \item \textbf{Comparaci\'on con m\'etodos alternativos}: Aplicar el \textit{pipeline} a
    modelos generados por RoseTTAFold \citep{baek2021accurate}, ESMFold
    \citep{lin2023evolutionary} y OpenFold para establecer un \textit{benchmark} comparativo
    con metodolog\'ia uniforme.

    \item \textbf{An\'alisis por familia estructural}: Estratificar por familia
    SCOP/CATH para identificar familias proteicas donde AlphaFold2 presenta mayor
    o menor precisi\'on, y correlacionar con la profundidad del MSA disponible.

    \item \textbf{Conformaciones m\'ultiples}: Evaluar la capacidad de AlphaFold2 para
    capturar ens\'amblajes conformacionales mediante \textit{dropout sampling}
    (subsampling de MSA durante la inferencia), t\'ecnica que algunos estudios
    recientes han propuesto para explorar el paisaje conformacional.

    \item \textbf{Aplicaci\'on a dise\~no de f\'armacos}: Validar la utilidad
    pr\'actica de los modelos AlphaFold2 en protocolos de docking molecular y
    \textit{virtual screening}, cuantificando el impacto del RMSD del modelo en
    las tasas de acierto (\textit{hit rates}) del cribado.
\end{enumerate}

% ═══════════════════════════════════════════════════════════════════════════
\subsection{Reflexi\'on final}
\label{subsec:reflexion_final}
% ═══════════════════════════════════════════════════════════════════════════

Este trabajo demuestra dos lecciones fundamentales. La primera es que AlphaFold2
constituye una herramienta de predicci\'on estructural extraordinariamente precisa:
con un RMSD mediano de \SI{1,26}{\angstrom} sobre 10.432 prote\'inas diversas, sus
predicciones se sit\'uan dentro del rango de variabilidad experimental, haciendo que
la estructura computacional sea indistinguible de la experimental para la mayor\'ia
de aplicaciones pr\'acticas.

La segunda lecci\'on es que \textbf{la metodolog\'ia de comparaci\'on es tan
importante como el algoritmo evaluado}. Un error en el mapeo de residuos puede
transformar un RMSD de \SI{0,64}{\angstrom} en \SI{22,54}{\angstrom}, convirtiendo
una herramienta excelente en aparentemente deficiente. La rigurosidad metodol\'ogica
---incluyendo el uso de SIFTS para la correspondencia correcta de posiciones, la
diversificaci\'on del \textit{dataset} para eliminar sesgos, y la cuantificaci\'on bayesiana
de la incertidumbre--- no es un detalle t\'ecnico sino un requisito fundamental para
cualquier evaluaci\'on estructural comparativa.

En un momento en que la comunidad cient\'ifica debate el papel de la predicci\'on
computacional frente a la determinaci\'on experimental de estructuras, estos
resultados aportan evidencia cuantitativa s\'olida: AlphaFold2, utilizado e
interpretado correctamente, produce modelos estructurales de calidad comparable a la
experimental. La transici\'on hacia AlphaFold3 ---con su capacidad para modelar
complejos multicomponente, \'acidos nucleicos e interacciones con ligandos--- promete
ampliar a\'un m\'as el impacto de la predicci\'on estructural computacional, y los
marcos metodol\'ogicos establecidos en este estudio ser\'an directamente aplicables
a su evaluaci\'on.
